% §6 Taxonomy of Abstraction Failure and Success Patterns
% Status: Written 2026-03-25 — full rewrite with validated E1–E7 evidence.

\section{Taxonomy of Abstraction Patterns}
\label{sec:taxonomy}

The experimental results of Section~\ref{sec:results} reveal that abstraction
failures and successes are not random --- they are \emph{structural}, recurring
whenever the same combination of workload characteristics, problem size, and
framework properties is present. This section formalizes these observations into
an empirical taxonomy of eight abstraction patterns (P001--P008): six failure
modes, one success pattern, and one unconfirmed hypothesis. Each pattern is
defined by a symptom signature, a root-cause category, quantitative evidence
from one or more experiments, and a recommended mitigation. The taxonomy
provides the vocabulary that the decision framework of
Section~\ref{sec:framework} operates on.

% ------------------------------------------------------------------
\subsection{Root-Cause Categories}
\label{subsec:categories}
% ------------------------------------------------------------------

Every pattern identified in this taxonomy is attributed to one of four
root-cause categories, each associated with a distinct diagnostic methodology:

\begin{enumerate}
    \item \textbf{Runtime Coordination Overhead} --- per-launch dispatch cost,
          fence and synchronization latency, or task scheduling overhead that
          the abstraction layer adds to each GPU kernel invocation. Detected
          via \texttt{nsys}/\texttt{rocprof} timeline analysis and CUDA/HIP
          event timing. Three of the five validated failure patterns (P001,
          P006, P008) fall in this category.

    \item \textbf{Abstraction API Limitation} --- the abstraction's programming
          interface cannot express the algorithm required for optimal
          performance (e.g., scoped shared-memory allocation for tiled DGEMM).
          Detected via source-code analysis comparing abstraction capabilities
          against native kernel requirements.

    \item \textbf{Memory Model Mismatch} --- the abstraction forces a
          non-optimal data layout or access pattern, causing coalescing
          degradation or unnecessary memory traffic. Detected via L2 cache hit
          rate and memory throughput counters. This category also encompasses
          the inverse case (P005) where the abstraction's default layout
          \emph{improves} coalescing relative to the native baseline.

    \item \textbf{Irregular Work Distribution} --- workload-inherent load
          imbalance (e.g., power-law degree distributions) is amplified by the
          abstraction's thread-to-work mapping, compounding with dispatch
          overhead. Detected via warp activity counters and per-row work
          statistics.
\end{enumerate}

\noindent A fifth category, \textbf{Compiler Backend Failure} (suboptimal code
generation by the abstraction's compiler backend), was hypothesized in our
initial experimental design. While suggestive evidence exists (P002), we were
unable to confirm this category independently due to hardware profiling
constraints (Section~\ref{subsec:unconfirmed}). It is not included as a
primary category pending further validation.

% ------------------------------------------------------------------
\subsection{Validated Failure Patterns}
\label{subsec:failure_patterns}
% ------------------------------------------------------------------

This subsection presents the five failure patterns for which we have
sufficient quantitative evidence to establish structural recurrence
across experiments.

% --- P001 ---
\paragraph{P001: Launch Overhead Dominance.}
\textit{Status: Validated.}
\textit{Category: Runtime Coordination Overhead.}

\noindent\textbf{Symptom.}
Abstraction efficiency drops to 0.17--0.70 when single-kernel execution time
$T_k$ is comparable to or smaller than the abstraction's per-launch overhead.

\noindent\textbf{Root cause.}
Each abstraction layer adds per-invocation overhead beyond the native
\texttt{cudaLaunchKernel} or \texttt{hipLaunchKernel} cost: runtime dispatch,
template instantiation, and synchronization logic that is not fully inlined.
When $T_k < \text{overhead}$, overhead dominates total execution time.
Per-launch overhead varies by framework: SYCL/HSA $\approx 19$\,\textmu{}s,
Julia/CUDA $\approx 7$\,\textmu{}s, Kokkos/CUDA $\approx 5$\,\textmu{}s,
RAJA/CUDA $\approx 2$\,\textmu{}s (Table~\ref{tab:overhead}).

\noindent\textbf{Canonical evidence.}
SYCL on E1 (STREAM Triad) at small $N$ on AMD MI300X: efficiency $= 0.167$
(527 vs.\ 3154\,GB/s native). The HSA dispatch path adds $\approx 19$\,\textmu{}s
per launch against a kernel time of $\approx 3.8$\,\textmu{}s, yielding a
predicted efficiency of $3.8 / (3.8 + 19) = 0.167$ --- an exact match with
observed performance (Section~\ref{subsec:e1}).

\noindent\textbf{Secondary evidence.}
Julia/E3/small/RTX~5060: efficiency $= 0.704$ ($T_k \approx 14$\,\textmu{}s,
overhead $\approx 7$\,\textmu{}s). All abstractions on E7 (N-Body) at small
$N$ (32K atoms): efficiency 0.39--0.67 across all four frameworks
(Section~\ref{subsec:e7}).

\noindent\textbf{PPC impact.} 0.17--0.70 depending on the $T_k$/overhead ratio.

\noindent\textbf{Mitigation.}
Use native implementations for kernels with $T_k < 50$\,\textmu{}s. Coarsen
task granularity (fuse multiple small kernels into a single launch) to
amortize overhead before adopting an abstraction.

\begin{table}[t]
\centering
\caption{Calibrated per-launch overhead by abstraction and backend.
Values are device-side effective overhead (\textmu{}s) after JIT warmup,
computed as $\text{overhead} = T_k \times (1 - e) / e$ from observed
efficiency $e$ across E1--E7 configurations.}
\label{tab:overhead}
\begin{tabular}{llcc}
\hline
\textbf{Abstraction} & \textbf{Backend} & \textbf{Median (\textmu{}s)}
& \textbf{Range (\textmu{}s)} \\
\hline
RAJA     & CUDA & 2  & 1--4  \\
RAJA     & HIP  & 4  & 3--5  \\
Kokkos   & CUDA & 5  & 4--7  \\
Kokkos   & HIP  & 4  & 3--5  \\
Julia    & CUDA & 7  & 5--10 \\
Julia    & ROCm & 4  & 3--6  \\
SYCL     & HSA  & 19 & 15--25 \\
\hline
\end{tabular}
\end{table}

% --- P004 ---
\paragraph{P004: API Expressivity Gap.}
\textit{Status: Validated.}
\textit{Category: Abstraction API Limitation.}

\noindent\textbf{Symptom.}
Size-independent efficiency degradation (PPC does \emph{not} improve with
increasing $N$) on compute-bound kernels that require shared-memory tiling.

\noindent\textbf{Root cause.}
RAJA's \texttt{forall} and Julia's \texttt{@cuda}/\texttt{@roc} macro
interfaces lack scoped shared-memory allocation primitives. Without
shared-memory tiling, DGEMM falls back to global-memory-only access with
$O(N^3)$ DRAM traffic versus $O(N^2)$ for the tiled native kernel. This is an
\emph{algorithmic} limitation, not a runtime cost: increasing $N$ does not
help because the memory-traffic asymptote is fundamentally different.

\noindent\textbf{Canonical evidence.}
RAJA on E2 (DGEMM) on AMD MI300X: efficiency $= 0.316$ (small) $\rightarrow
0.236$ (medium) $\rightarrow 0.189$ (large) --- the \emph{worsening} trend
with $N$ confirms an algorithmic root cause (Section~\ref{subsec:e2}). On the
same kernel, Kokkos \texttt{TeamPolicy} achieves 0.963 and SYCL
\texttt{local\_accessor} achieves 0.972, isolating API expressivity as the
discriminating factor.

\noindent\textbf{PPC impact.} 0.19--0.35 for affected abstractions (RAJA, Julia)
on compute-bound tiled kernels.

\noindent\textbf{Mitigation.}
Use Kokkos \texttt{TeamPolicy} with scratch-pad memory, SYCL
\texttt{local\_accessor}, or vendor-tuned libraries (cuBLAS, rocBLAS) for
kernels that require shared-memory tiling.

% --- P006 ---
\paragraph{P006: Tiling Policy Overhead.}
\textit{Status: Partially validated.}
\textit{Category: Runtime Coordination Overhead.}

\noindent\textbf{Symptom.}
Kokkos \texttt{MDRangePolicy} is 15--25\% slower than flat-loop dispatch at
small $N$; shared-memory tiling is counterproductive for data structures
without spatial locality.

\noindent\textbf{Root cause.}
\texttt{MDRangePolicy<Rank<3>{>}} with tile hints generates boundary-masking
conditionals at tile edges. At small $N$ (e.g., $32^3$), boundary overhead
constitutes 20--30\% of warp execution cycles. Additionally, tiling assumes
spatial data reuse; when the data structure does not provide it (e.g., CSR
neighbor lists in N-body), shared-memory staging adds pressure and barrier
cost without compensating reuse.

\noindent\textbf{Evidence.}
E3 (3D Stencil) on RTX~5060: Kokkos efficiency $= 0.760$ versus RAJA
$= 0.942$ with identical kernel logic but different dispatch policies
(Section~\ref{subsec:e3}). E7 (N-Body): native tiled kernel is 15--34\%
\emph{slower} than native untiled at all problem sizes for CSR-based force
calculation (Section~\ref{subsec:e7}).

\noindent\textbf{PPC impact.} 0.65--0.85 at small $N$; recovers to $\geq 0.90$ at
large $N$ as boundary overhead amortizes.

\noindent\textbf{Mitigation.}
Reduce tile size proportionally to $N$; switch to \texttt{RangePolicy} with
manual index arithmetic for small grids. Profile tiling benefit before
assuming it helps --- tiling is data-structure-dependent, not just
architecture-dependent.

% --- P007 ---
\paragraph{P007: Load Imbalance Amplification.}
\textit{Status: Validated.}
\textit{Category: Irregular Work Distribution.}

\noindent\textbf{Symptom.}
Efficiency collapses \emph{below} the level predicted by P001 alone ---
launch overhead and load imbalance interact multiplicatively.

\noindent\textbf{Root cause.}
Power-law degree distributions ($\sigma = 2.5$) create extreme row-length
variance. One-thread-per-row mapping leaves most warps stalled waiting on the
few long-row threads. When P001 is already active (small $N$, high overhead),
the imbalance compounds: idle warp cycles increase the \emph{effective}
overhead fraction per launch.

\noindent\textbf{Canonical evidence.}
Julia on E4 (SpMV) at small $N$ on RTX~5060: power-law efficiency $= 0.449$
versus Laplacian efficiency $= 0.561$ --- same kernel, same $N$, but
$\sigma = 2.5$ adds a 0.11 multiplicative penalty beyond P001
(Section~\ref{subsec:e4}). The compound model
$\text{PPC} = \text{PPC}_{\text{P001}} \times (1 - 0.10 \times
\min(\sigma, 3.0))$ predicts 0.42; observed 0.449.

\noindent\textbf{PPC impact.} 0.40--0.70 at small $N$ with power-law inputs;
near-native at large $N$ where both P001 and P007 amortize.

\noindent\textbf{Mitigation.}
Warp-per-row CSR (multiple threads per row), segmented reductions
(\texttt{DeviceSegmentedReduce}), or row reordering to group similar-length
rows within each warp. Native implementations recommended for irregular
workloads at small $N$.

% --- P008 ---
\paragraph{P008: Level-Set Dispatch Amplification.}
\textit{Status: Validated --- strongest quantitative finding.}
\textit{Category: Runtime Coordination Overhead.}

\noindent\textbf{Symptom.}
Efficiency decays with the number of dependency levels $L$, \emph{not} with
problem size $N$. Level-set algorithms accumulate
$O(L \times \text{overhead})$ total penalty.

\noindent\textbf{Root cause.}
Level-set algorithms (SpTRSV forward substitution, BFS frontier expansion)
require one GPU kernel launch per dependency level to preserve correctness.
Each launch carries the abstraction's per-invocation overhead
(Table~\ref{tab:overhead}). Total overhead $= L \times \text{overhead}$,
which is workload-structure-dependent. The distinguishing test: fix $N$, vary
$L$ --- if efficiency decays with $L$ but not $N$, the pattern is P008.

\noindent\textbf{Canonical evidence.}
SYCL on E5 (SpTRSV) with Laplacian matrix ($L = 181$) on AMD MI300X:
predicted PPC $= T_{\text{compute}} / (T_{\text{compute}} + 181 \times
0.070\,\text{ms}) = 0.490$; observed PPC $= 0.490$ --- an exact model match
(Section~\ref{subsec:e5}). Julia on E5/Laplacian/large on RTX~5060: efficiency
$= 0.464$ ($181 \times 25$\,\textmu{}s $= 4.5$\,ms accumulated overhead).
Julia on E6 (BFS): efficiency 5--6\% due to multi-kernel dispatch per level
(Section~\ref{subsec:e6}).

\noindent\textbf{Inverse pattern.}
When levels are \emph{wide} (many rows per level), P008 reverses: RAJA on
E5/random/large on RTX~5060 achieves efficiency $= 1.299$, and Kokkos
$= 1.113$. Wide levels make per-launch overhead negligible relative to
compute, and RAJA/Kokkos fence implementations pipeline more efficiently
than native \texttt{cudaDeviceSynchronize}.

\noindent\textbf{SYCL diagnostic axis.}
P008 explains 90\% of SYCL's performance variance across experiments:
single-kernel workloads (E3: 0.987, E7: 0.92) deliver excellent efficiency,
while multi-kernel workloads (E5: 0.49, E6: 0.30) collapse. The
discriminating variable is not architecture but workload structure ---
specifically, whether the algorithm requires $L > 1$ launches per invocation.

\noindent\textbf{PPC impact.} Julia: 0.35--0.46; SYCL: 0.49--0.54;
RAJA/Kokkos: 0.80--1.30 depending on level width.

\noindent\textbf{Mitigation.}
Prefer RAJA or Kokkos for level-set algorithms (per-launch overhead 2--5\,\textmu{}s
versus Julia 7--25\,\textmu{}s or SYCL 19\,\textmu{}s). Fuse independent
levels where correctness permits. Pre-compute level structure and batch
independent rows into coarsened super-levels.

% ------------------------------------------------------------------
\subsection{Success Patterns and Unconfirmed Hypotheses}
\label{subsec:success_patterns}
% ------------------------------------------------------------------

% --- P005 ---
\paragraph{P005: Layout-Induced Coalescing Advantage.}
\textit{Status: Validated --- success pattern.}
\textit{Category: Memory Model Mismatch (inverse).}

\noindent\textbf{Symptom.}
Abstraction efficiency exceeds 1.0 --- the abstraction \emph{outperforms}
the native baseline.

\noindent\textbf{Root cause.}
Julia's column-major \texttt{CuArray}s map warp threads
(\texttt{threadIdx.x} $\rightarrow$ row index) to consecutive elements of the
same matrix column, achieving perfect coalescing on the $A$ matrix and
broadcast access on $B$. The native tiled C kernel uses row-major storage and
incurs 512 \texttt{\_\_syncthreads()} barriers at $N = 8192$ (256 tiles
$\times$ 2 barriers per tile). At arithmetic intensity $\approx 683$
FLOP/byte, synchronization overhead exceeds shared-memory bandwidth savings.

\noindent\textbf{Evidence.}
Julia on E2 (DGEMM) on RTX~5060: efficiency $= 1.25\times$ native at large
$N$ (Section~\ref{subsec:e2}). Fresh-state profiling confirms the advantage
is size-independent: $1.29\times$ at medium $N$, $1.20\times$ at large $N$.

\noindent\textbf{Platform asymmetry.}
P005 is confirmed on NVIDIA RTX~5060 \emph{only}. On AMD MI300X, Julia naive
DGEMM achieves efficiency $= 0.306$ --- the column-major layout advantage
is insufficient when the native AMD baseline is $4\times$ higher performing.
This platform-generation reversal (Godoy et al.~\cite{Godoy2023} found
$0.90$+ on CDNA2/MI250X) underscores that abstraction outcomes are
architecture-specific.

\noindent\textbf{Implication.}
Layout conventions can dominate algorithmic overhead in compute-bound regimes.
The expert assumption that ``tiling always improves performance'' fails when
synchronization cost exceeds the memory-traffic reduction from shared-memory
staging.

% --- P002/P003 ---
\label{subsec:unconfirmed}
\paragraph{Unconfirmed patterns.}
Two patterns from our initial hypothesis set were not independently confirmed:

\textit{P002: Compiler Backend Unrolling Failure} (partially validated).
The original hypothesis attributed E2 RAJA naive degradation on NVIDIA to
suboptimal NVCC code generation for template-instantiated loops. Subsequent
analysis reclassified this as P004 (API Expressivity Gap): RAJA's performance
deficit is size-independent and persists identically on AMD, ruling out an
NVIDIA-specific compiler cause. Independent confirmation requires
\texttt{ncu} hardware counter analysis, which is currently blocked by
\texttt{RmProfilingAdminOnly=1} on the RTX~5060 driver. P002 remains
plausible as a contributing factor but cannot be isolated from P004 with
available data.

\textit{P003: Memory Layout Mismatch} (hypothesis).
E3 Kokkos small-$N$ efficiency ($0.760$) was initially attributed to a layout
mismatch, but RAJA achieves $0.942$ on the same kernel with identical data
layout --- the gap is better explained by P006 (\texttt{MDRangePolicy}
overhead). No independent P003 evidence was found across E1--E7. The pattern
may manifest on future hardware or with workloads not covered by our
experimental matrix.

% ------------------------------------------------------------------
\subsection{Taxonomy Summary}
\label{subsec:taxonomy_summary}
% ------------------------------------------------------------------

Table~\ref{tab:taxonomy_summary} presents all eight patterns with their
validation status, root-cause category, PPC impact range, and primary
experimental evidence.

\begin{table*}[t]
\centering
\caption{Empirical taxonomy of abstraction patterns. Status indicates
the level of experimental evidence: \textit{validated} (structural
recurrence confirmed across multiple configurations), \textit{partial}
(suggestive evidence, pending further instrumentation), \textit{hypothesis}
(no independent confirmation).}
\label{tab:taxonomy_summary}
\begin{tabular}{llllcl}
\hline
\textbf{ID} & \textbf{Pattern Name} & \textbf{Type} & \textbf{Category}
& \textbf{PPC Range} & \textbf{Primary Evidence} \\
\hline
P001 & Launch Overhead Dominance      & Failure  & Runtime Overhead
     & 0.17--0.70 & E1 SYCL/small/AMD \\
P002 & Compiler Backend Failure       & Failure  & Compiler Backend
     & 0.25--0.50 & E2 (inconclusive) \\
P003 & Memory Layout Mismatch         & Failure  & Memory Model
     & 0.60--0.80 & --- \\
P004 & API Expressivity Gap           & Failure  & API Limitation
     & 0.19--0.35 & E2 RAJA/all sizes/AMD \\
P005 & Layout-Induced Coalescing      & Success  & Memory Model
     & $> 1.0$    & E2 Julia/large/RTX \\
P006 & Tiling Policy Overhead         & Failure  & Runtime Overhead
     & 0.65--0.85 & E3 Kokkos/small/RTX; E7 tile \\
P007 & Load Imbalance Amplification   & Failure  & Irregular Work
     & 0.40--0.70 & E4 Julia/power\_law/small \\
P008 & Level-Set Dispatch Amplification & Failure & Runtime Overhead
     & 0.35--1.30 & E5 SYCL/laplacian/AMD \\
\hline
\end{tabular}
\end{table*}

% ------------------------------------------------------------------
\subsection{Cross-Reference and Discussion}
\label{subsec:taxonomy_map}
% ------------------------------------------------------------------

Table~\ref{tab:taxonomy_crossref} maps each pattern to the experiments in
which it was observed, revealing three structural properties of the taxonomy.

\begin{table}[t]
\centering
\caption{Pattern--experiment cross-reference. $\bullet$ = primary evidence;
$\circ$ = secondary evidence. E1: STREAM Triad; E2: DGEMM; E3: 3D Stencil;
E4: SpMV; E5: SpTRSV; E6: BFS; E7: N-Body.}
\label{tab:taxonomy_crossref}
\begin{tabular}{lccccccc}
\hline
\textbf{Pattern} & \textbf{E1} & \textbf{E2} & \textbf{E3} & \textbf{E4}
& \textbf{E5} & \textbf{E6} & \textbf{E7} \\
\hline
P001 & $\bullet$ &           & $\circ$   & $\circ$   &           &           & $\circ$   \\
P002 &           & $\circ$   &           &           &           &           &           \\
P003 &           &           &           &           &           &           &           \\
P004 &           & $\bullet$ & $\circ$   &           &           &           &           \\
P005 &           & $\bullet$ &           &           &           &           &           \\
P006 &           &           & $\bullet$ &           &           &           & $\bullet$ \\
P007 &           &           &           & $\bullet$ &           &           &           \\
P008 &           &           &           &           & $\bullet$ & $\circ$   & $\circ$   \\
\hline
\end{tabular}
\end{table}

\paragraph{Concentration.}
Three patterns --- P001 (Launch Overhead), P004 (API Expressivity), and P008
(Level-Set Dispatch) --- account for approximately 80\% of all observed
failure configurations with PPC $< 0.60$. This concentration suggests that
practical abstraction selection can be driven by a small number of diagnostic
checks rather than exhaustive profiling.

\paragraph{Shared mechanisms at different scales.}
P001 and P008 share the same root cause (per-launch dispatch overhead) but
operate at different scales: P001 manifests when a \emph{single} kernel's
execution time is comparable to overhead; P008 manifests when
\emph{accumulated} overhead across $L$ levels dominates total runtime. The
distinction is practically important: P001 can be resolved by coarsening a
single kernel, while P008 requires algorithmic changes (level fusion) or
framework selection (RAJA/Kokkos over Julia/SYCL for level-set workloads).

\paragraph{Multiplicative co-occurrence.}
When multiple patterns are simultaneously active, their PPC effects compound
multiplicatively rather than additively. P007 (Load Imbalance) compounds with
P001 (Launch Overhead) at small $N$ on power-law inputs: P001 alone predicts
efficiency $\approx 0.56$ for Julia/E4/small, but the P001$\times$P007
compound yields 0.449. This multiplicative interaction means that developers
facing irregular workloads at small problem sizes must address \emph{both}
dispatch overhead and load imbalance to recover acceptable performance.

\paragraph{Structural recurrence.}
All five validated patterns are structural --- they recur whenever the same
combination of workload characteristics and overhead profile is present,
independent of specific framework version or driver release. P001 appears
identically in E1 (STREAM), E3 (Stencil), E4 (SpMV), and E7 (N-Body)
whenever $T_k < \text{overhead}$. P008 appears in both E5 (SpTRSV) and E6
(BFS) whenever $L$ is high, and its \emph{absence} in E7 (single-kernel
N-Body) confirms the $L$-dependence. This structural property is what makes
the patterns predictive rather than merely descriptive.

\paragraph{Limitations.}
Two of the eight hypothesized patterns (P002, P003) could not be confirmed
due to hardware profiling constraints. The RTX~5060's
\texttt{RmProfilingAdminOnly=1} driver restriction blocks \texttt{ncu}
hardware performance counters required for compiler code-generation analysis
(P002) and L2 cache hit-rate measurement (P003). Future work should validate
these patterns on hardware with full performance monitoring unit access.

\medskip
\noindent These patterns and their quantitative signatures form the input
vocabulary for the five calibrated decision rules presented in
Section~\ref{sec:framework}.
